\subsection{Yêu cầu phi chức năng}

Các yêu cầu phi chức năng nhằm đảm bảo khả năng vận hành, độ tin cậy và an toàn thông tin:

\begin{enumerate}
    \item Hiệu năng và khả năng mở rộng (Performance \& Scalability):
    \begin{itemize}[label={-}]
        \item Xử lý đồng thời: Sử dụng kiến trúc I/O đa kênh (I/O Multiplexing) với epoll và Thread Pool nhằm duy trì luồng điều phối phi đồng bộ (Non-blocking). Giải quyết bài toán C10K với độ trễ thấp.
        \item Tối ưu băng thông: Áp dụng giao thức nhị phân (Custom Binary Protocol) để tối thiểu hóa kích thước gói tin và hao phí phân tích cú pháp (Parsing).
    \end{itemize}

    \item Bảo mật và toàn vẹn dữ liệu (Security \& Integrity):
    \begin{itemize}[label={-}]
        \item Mã hóa kênh truyền: Mã hóa tải trọng bằng AES-256-GCM, trao đổi khóa phiên (Session Key) qua RSA-2048.
        \item Toàn vẹn thông điệp: Áp dụng kiểm lỗi CRC32 và mã xác thực HMAC cho từng gói tin.
        \item An toàn lưu trữ: Băm mật khẩu một chiều trước khi lưu trữ. Sử dụng Write-Ahead Logging (WAL) cho cơ sở dữ liệu SQLite để duy trì tính nhất quán.
    \end{itemize}

    \item Tính tin cậy và khả năng chịu lỗi (Reliability \& Fault Tolerance):
    \begin{itemize}[label={-}]
        \item Quản lý tài nguyên: Tự động thu hồi tài nguyên (socket, phiên đăng nhập, vùng nhớ) của kết nối vi phạm thời gian chờ (Timeout) hoặc ngắt đột ngột, ngăn chặn cạn kiệt mô tả tệp (File Descriptor Exhaustion) và rò rỉ bộ nhớ (Memory Leak).
        \item Xử lý luồng mạng: Triển khai bộ đệm (Network Buffer) độc lập cho từng kết nối nhằm xử lý hiện tượng phân mảnh (TCP Fragmentation) và dính gói (TCP Sticky Packets).
    \end{itemize}
\end{enumerate}